\chapter{CƠ SỞ KỸ THUẬT VÀ CÔNG NGHỆ NỀN TẢNG}

\section{Phương pháp luận và tiêu chuẩn kỹ thuật áp dụng}

Để đảm bảo tính chuẩn mực học thuật và độ tin cậy kỹ thuật trong công nghiệp, hệ thống MarketFlow AI được thiết kế và hiện thực hóa dựa trên khung phương pháp luận Nghiên cứu Khoa học Thiết kế (\textit{Design Science Research Methodology} -- DSRM) đề xuất bởi Peffers và cộng sự \cite{peffers2007dsrm}. Quá trình phát triển tuân thủ một chu trình lặp chặt chẽ: xuất phát từ việc nhận diện điểm nghẽn thực tế trong vận hành tiếp thị và tích hợp AI tác tử (khảo sát bộ nhớ P0 Spike), xác lập các mục tiêu giải pháp khả thi, kiến tạo cấu trúc phần mềm thực thi, và tiến hành kiểm định tự động đa tầng.

Hệ thống tích hợp đồng bộ các tiêu chuẩn kỹ thuật quốc tế:
\begin{itemize}
 \item \textbf{ISO/IEC/IEEE 29148:2018 (Requirements Engineering)}: Toàn bộ 14 yêu cầu chức năng (FR01--FR14) và 6 yêu cầu phi chức năng (NFR01--NFR06) được định danh duy nhất, phân rã nguyên tử, có điều kiện nghiệm thu rõ ràng và liên kết trực tiếp với 151 ca kiểm thử tự động trong Ma trận truy xuất nguồn gốc (RTM).
 \item \textbf{ISO/IEC/IEEE 42010:2011 (Architecture Description)}: Kiến trúc hệ thống được phân tách theo góc nhìn của các bên liên quan, mô hình hóa theo khung chuẩn C4 (Context, Container, Component) và tách biệt rõ ranh giới giữa dịch vụ nghiệp vụ (\texttt{backend}) với dịch vụ thực thi AI (\texttt{ai-runner}).
 \item \textbf{ISO/IEC 25010:2011 (System and Software Quality Models)}: Định chuẩn các thuộc tính chất lượng phần mềm bao gồm: Tính đúng đắn chức năng, Hiệu năng phản hồi, Tính khả dụng, Độ tin cậy (chịu lỗi khi dịch vụ AI ngoại tuyến), và An toàn bảo mật.
 \item \textbf{NIST AI Risk Management Framework (AI RMF 1.0)}: Khung quản lý rủi ro trí tuệ nhân tạo của Viện Tiêu chuẩn và Công nghệ Quốc gia Hoa Kỳ \cite{nist-airmf,nist-genai-profile}, đặc biệt là nguyên tắc kiểm soát ảo giác bằng truy vấn công cụ (\textit{tool-grounded generation}), cô lập quyền thực thi qua mã thông báo tạm thời (\texttt{Run-Token} JWT) và bảo đảm sự phê duyệt của con người (\textit{Proposal-Only HITL}).
\end{itemize}

\section{Khảo sát và phân tích lựa chọn công nghệ cốt lõi}

Hệ thống MarketFlow AI được xây dựng trên một ngăn xếp công nghệ hiện đại, cân bằng tối ưu giữa hiệu năng thực thi cao, độ an toàn khi điều phối tác tử AI và tính linh hoạt khi triển khai trên cả môi trường cục bộ lẫn máy chủ đám mây. Bảng~\ref{tab:tech-stack} tổng hợp các công nghệ nền tảng đã được triển khai trong thực tế.

\begin{table}[!htbp]
\centering
\small
\caption{Ngăn xếp công nghệ hiện thực hóa hệ thống MarketFlow AI}
\label{tab:tech-stack}
\begin{tabularx}{\textwidth}{|L{2.8cm}|L{3.2cm}|Y|}
\hline
\thead{Thành phần} & \thead{Công nghệ lựa chọn} & \thead{Phân tích kỹ thuật và lý do lựa chọn thực tế}\\
\hline
Backend Web API \& MCP Server & FastAPI (Python 3.12+), Pydantic v2, FastMCP & Framework ASGI hiệu năng cao; tích hợp Pydantic v2 kiểm định dữ liệu tự động; nhúng trực tiếp máy chủ giao thức \texttt{FastMCP} (\texttt{stateless\_http=True}) tại đường dẫn nội bộ \texttt{/mcp} cung cấp 14 công cụ nghiệp vụ.\\
\hline
AI Runner \& Agent Runtime & FastAPI (\texttt{:8100}), OpenCode CLI v1.19.4 & Dịch vụ vi mô tách biệt quản lý hàng đợi FIFO (\texttt{RUNNER\_CONCURRENCY=1}), khởi chạy tiến trình con \texttt{opencode run --agent marketflow --format json} trong môi trường cô lập (tài khoản \texttt{mfai}, giới hạn \texttt{MemoryMax=896M}).\\
\hline
Hệ sinh thái Mô hình LLM & 13 mô hình trên 5 nền tảng (OpenCode, Groq, OpenRouter, Gemini) & Mô hình chính \texttt{opencode/deepseek-v4.1-flash-free} kết hợp 12 mô hình dự phòng (GLM-5.3, MiMo-v2.6, GPT-OSS-120B, Qwen3.8-27B, Nemotron-3.5, Gemini 3.5/3.1 Flash Lite) giúp duy trì chi phí vận hành 0 USD và chống đứt gãy dịch vụ.\\
\hline
Frontend Client & React 18, TypeScript, Vite, Tailwind CSS v4 & Kiến trúc SPA phân tách theo miền nghiệp vụ (\texttt{src/features/*}); tích hợp TanStack Query v5, ECharts (biểu đồ KPI), FullCalendar (lịch xuất bản) và luồng đọc sự kiện thời gian thực \texttt{fetch} + \texttt{ReadableStream} (SSE).\\
\hline
Cơ sở dữ liệu & PostgreSQL 16 (Production) \& SQLite WAL (Dev/Test) & Quản lý 14 bảng quan hệ qua SQLAlchemy 2.0 ORM và Alembic migrations; sử dụng PostgreSQL 16 trên VPS đảm bảo tính bền vững đồng thời cao, kết hợp SQLite (WAL mode + \texttt{foreign\_keys=ON}) cho kiểm thử nhanh.\\
\hline
Đóng gói, Hạ tầng \& Kiểm thử & Docker Compose, Systemd, Nginx, Pytest, Playwright & Đóng gói 3 container (\texttt{postgres}, \texttt{backend}, \texttt{ai-runner}); vận hành thực tế qua Cloudflare Tunnel + Nginx + Systemd; kiểm chứng bằng 151 ca kiểm thử tự động (Pytest + Playwright E2E).\\
\hline
\end{tabularx}
\end{table}

\subsection{Backend Framework: FastAPI, Pydantic v2 và FastMCP nội bộ}
FastAPI được lựa chọn làm máy chủ trung tâm nhờ các ưu điểm kỹ thuật vượt trội:
\begin{enumerate}
 \item \textbf{Kiến trúc đồng bộ kết hợp truyền phát SSE}: Các điểm cuối REST API nghiệp vụ (\texttt{/api/v1/*}) được xử lý rõ ràng, trong khi điểm cuối hội thoại AI (\texttt{POST /api/v1/assistant/threads/\{id\}/messages}) sử dụng \texttt{StreamingResponse} chuẩn \texttt{text/event-stream} để đẩy từng sự kiện \texttt{meta}, \texttt{status}, \texttt{tool}, \texttt{delta}, \texttt{action} và \texttt{done} về trình duyệt theo thời gian thực.
 \item \textbf{Kiểm định dữ liệu chặt chẽ qua Pydantic v2}: Hệ thống tận dụng Pydantic v2 để kiểm tra tính hợp lệ của ngày tháng (\texttt{start\_date} $\le$ \texttt{end\_date}), ngân sách dương (\texttt{budget} $> 0$), số liệu phễu tiếp thị (\texttt{views} $\ge$ \texttt{clicks} $\ge$ \texttt{conversions}) và cấu trúc đầu vào của toàn bộ 14 công cụ MCP.
 \item \textbf{Nhúng trực tiếp máy chủ FastMCP}: Thay vì mở thêm một cổng mạng riêng biệt khó kiểm soát, mô-đun \texttt{backend/app/assistant/mcp\_server.py} gắn trực tiếp ứng dụng \texttt{FastMCP("marketflow-assistant", stateless\_http=True)} vào đường dẫn \texttt{/mcp} của FastAPI. Máy chủ Nginx chặn hoàn toàn truy cập \texttt{/mcp} từ bên ngoài (\texttt{location \^{}\~{} /mcp \{ return 404; \}}), chỉ cho phép tiến trình cục bộ trên cùng máy chủ gọi vào kèm mã thông báo xác thực.
\end{enumerate}

\subsection{Dịch vụ AI Runner, OpenCode CLI và Hệ sinh thái 13 mô hình đa nhà cung cấp}
Đáp ứng tiêu chí tích hợp AI hiện đại, có kiểm soát và tối ưu kinh tế của học phần AIA331, hệ thống thiết kế kiến trúc điều phối tác tử 2 tầng:
\begin{enumerate}
 \item \textbf{Cô lập tiến trình thực thi qua \texttt{ai-runner}}: Qua khảo sát thực nghiệm P0 Spike trên máy chủ mục tiêu, mỗi tiến trình \texttt{opencode run} (vận hành trên runtime Bun) tiêu thụ trung bình \texttt{\~{}490 MB RSS}. Để bảo vệ máy chủ VPS (2,5 GB RAM) khỏi nguy cơ tràn bộ nhớ (OOM), dịch vụ \texttt{ai-runner} (\texttt{ai-runner/app/runner.py}) sử dụng khóa điều phối \texttt{asyncio.Semaphore(1)} (\texttt{RUNNER\_CONCURRENCY=1}) để xếp hàng tuần tự các lượt gọi AI và chạy dưới tài khoản hệ thống \texttt{mfai} bị giới hạn cứng \texttt{MemoryMax=896M}.
 \item \textbf{Xác thực uỷ quyền ngắn hạn (\texttt{Run-Token} JWT)}: Khi người dùng gửi tin nhắn, Backend sinh một mã thông báo tạm thời (\texttt{typ="assistant\_run"}, thời hạn 10 phút) chứa định danh \texttt{user\_id}, \texttt{role} và \texttt{thread\_id}, truyền qua biến môi trường \texttt{MARKETFLOW\_RUN\_TOKEN} cho tiến trình \texttt{opencode}. Nhờ đó, mọi công cụ MCP mà AI gọi đều chịu sự kiểm soát phân quyền giống hệt người dùng đang đăng nhập, đồng thời không cần lưu trữ khóa phiên dài hạn trong tiến trình AI.
 \item \textbf{Hỗ trợ 13 mô hình LLM đa nhà cung cấp}: Giao diện \texttt{ModelSelector.tsx} và cấu hình \texttt{deploy/opencode/opencode.json} cho phép người dùng lựa chọn hoặc tự động chuyển đổi dự phòng (failover) giữa 13 mô hình thuộc 5 nền tảng: OpenCode Free (\texttt{deepseek-v4.1-flash-free}, \texttt{glm-5.3-flash-free}, \texttt{mimo-v2.6-flash-free}), Groq (\texttt{gpt-oss-120b}, \texttt{qwen3.8-27b}, \texttt{gpt-oss-20b}), OpenRouter (\texttt{nemotron-3.5-super-vl:free}, \texttt{nemotron-3-ultra-550b-a55b:free}, \texttt{nemotron-3.5-lightning:free}, \texttt{inkling:free}, \texttt{inkling-small:free}) và Google AI Studio (\texttt{gemini-3.5-flash-lite}, \texttt{gemini-3.1-flash-lite}).
\end{enumerate}

\subsection{Frontend: React 18, TypeScript, TanStack Query và ECharts}
Giao diện người dùng được xây dựng theo kiến trúc mô-đun hóa theo miền nghiệp vụ (\texttt{frontend/src/features/*}):
\begin{enumerate}
 \item \textbf{Quản lý trạng thái bất đồng bộ với TanStack Query v5}: Tự động đồng bộ bộ nhớ đệm (cache invalidation) ngay sau khi người dùng phê duyệt nội dung, nhập chỉ số CSV hoặc bấm \textit{Áp dụng} một thẻ đề xuất \texttt{ai\_action} từ khung chat Trợ lý AI.
 \item \textbf{Trực quan hóa dữ liệu chuyên sâu}: Sử dụng thư viện \texttt{ECharts} (\texttt{KpiChart.tsx}, \texttt{ChannelBreakdown.tsx}) để biểu diễn xu hướng lượt xem, lượt nhấp, chi phí và tỷ lệ giải ngân ngân sách; kết hợp \texttt{FullCalendar} (\texttt{CalendarPage.tsx}) hỗ trợ kéo-thả lịch xuất bản trực quan.
 \item \textbf{Khung Trợ lý AI tương tác thời gian thực (\texttt{AssistantDrawer.tsx} \& \texttt{AssistantPage.tsx})}: Hiển thị trực tiếp tiến trình suy luận, các thẻ gọi công cụ (\texttt{ToolCallChip.tsx}) và thẻ hành động đề xuất (\texttt{ProposedActionCard.tsx}) với hai nút quyết định rõ ràng: \textit{Áp dụng} và \textit{Bỏ qua}.
\end{enumerate}

\subsection{Cơ sở dữ liệu: PostgreSQL 16 (Production) và SQLite WAL (Dev/Test)}
Hệ thống sử dụng SQLAlchemy 2.0 ORM (\texttt{backend/app/models/}) kết hợp công cụ quản lý phiên bản lược đồ Alembic (\texttt{backend/alembic/}):
\begin{enumerate}
 \item \textbf{Mô hình dữ liệu chuẩn hóa 14 bảng}: Bao quát trọn vẹn 4 phân hệ: Quản trị \& Thương hiệu (\texttt{users}, \texttt{channels}, \texttt{brand\_profile}, \texttt{banned\_keywords}), Chiến dịch \& Phê duyệt HITL (\texttt{campaigns}, \texttt{campaign\_channels}, \texttt{contents}, \texttt{content\_reviews}), Vận hành \& Chỉ số (\texttt{schedules}, \texttt{metrics}, \texttt{notifications}) và Trợ lý AI (\texttt{ai\_threads}, \texttt{ai\_messages}, \texttt{ai\_actions}).
 \item \textbf{Tính tương thích kép PostgreSQL / SQLite}: Trên môi trường Production VPS và Docker Compose, hệ thống vận hành trên PostgreSQL 16 đảm bảo tính toàn vẹn giao dịch cao. Trong môi trường phát triển và chạy kiểm thử tự động (\texttt{pytest}), hệ thống hỗ trợ SQLite với chế độ \texttt{PRAGMA journal\_mode = WAL} và \texttt{PRAGMA foreign\_keys = ON} được kích hoạt tự động tại \texttt{backend/app/core/database.py}.
\end{enumerate}

\section{Kết luận chương}

Việc lựa chọn và kết hợp hài hòa giữa FastAPI, FastMCP, dịch vụ cô lập \texttt{ai-runner} với OpenCode CLI (13 mô hình LLM), React 18 + TypeScript và hệ quản trị CSDL PostgreSQL 16 / SQLite WAL đã tạo nên một nền tảng kỹ thuật hiện đại, an toàn và có tính thực tiễn cao cho MarketFlow AI. Chương tiếp theo sẽ đặc tả chi tiết các yêu cầu chức năng, phi chức năng, ma trận phân quyền RBAC và máy trạng thái duyệt nội dung HITL.
